%=============================================================================
% Engineering Design: km-Scale Wireless Power Demonstration
% via ψ-Channel Scalar Refractive Transmission
%
% Patent #12 — Provisional 63/865,283 (Aug 16 2025)
% Complete System Design Document
%=============================================================================

\documentclass[11pt,letterpaper]{article}

\usepackage[margin=1in]{geometry}
\usepackage{amsmath,amssymb}
\usepackage{graphicx}
\usepackage{booktabs}
\usepackage{siunitx}
\usepackage{hyperref}
\usepackage{xcolor}
\usepackage{tcolorbox}
\usepackage{enumitem}
\usepackage{float}
\usepackage{longtable}
\usepackage{fancyhdr}

\pagestyle{fancy}
\fancyhf{}
\fancyhead[L]{\small DFD Patent \#12 --- km-Scale Demonstration Design}
\fancyhead[R]{\small CONFIDENTIAL}
\fancyfoot[C]{\thepage}

\title{Engineering Design Document:\\
\textbf{km-Scale Wireless Power Transmission via\\
$\psi$-Channel Scalar Refractive Energy Delivery}\\[1em]
\large Phase II Demonstration System}

\author{DFD Research Programme\\
Provisional Patent 63/865,283}
\date{March 2026}

\begin{document}
\maketitle
\tableofcontents
\newpage

%=============================================================================
\section{Executive Summary}
%=============================================================================

This document presents the complete engineering design for a 1~km
wireless power transmission demonstration using the $\psi$-channel
scalar refractive mechanism predicted by Density Field Dynamics (DFD).
The system transmits energy by modulating the local $\psi$-field via
high-$Q$ electromagnetic cavities, propagating the $\psi$-perturbation
through a guided scalar waveguide, and reconverting to electromagnetic
energy at the receiver.

\begin{tcolorbox}[colback=blue!5!white, colframe=blue!50!black,
title=Key Design Parameters]
\begin{itemize}[nosep]
\item \textbf{Transmission distance}: 1.0~km
\item \textbf{Operating frequency}: $\Omega_\psi = 2.6$~GHz ($= 2 \times 1.3$~GHz)
\item \textbf{Transmitter array}: $10 \times 10$ phase-locked SRF cavities
\item \textbf{Receiver array}: $10 \times 10$ matched SRF cavities
\item \textbf{Channel maintenance}: 10 pump stations at 100~m spacing
\item \textbf{Target end-to-end efficiency}: $> 0.1\%$ (proof of concept)
\item \textbf{Total facility power}: 500~kW
\item \textbf{Estimated cost}: \$25--50M
\item \textbf{Timeline}: 36 months from funding to first power transmission
\end{itemize}
\end{tcolorbox}

%=============================================================================
\section{System Overview}
%=============================================================================

\subsection{Functional Architecture}

The demonstration system consists of five major subsystems:

\begin{enumerate}
\item \textbf{Transmitter Station (TX)}: Converts grid electrical power
to $\psi$-mode radiation via a phased SRF cavity array.
\item \textbf{Guided $\psi$-Channel}: A linear chain of pump stations
maintaining the background $\psi$-gradient waveguide.
\item \textbf{Receiver Station (RX)}: Converts $\psi$-mode energy back
to EM and rectifies to DC power.
\item \textbf{Diagnostic Network}: Distributed sensors measuring
$\psi$-field, EM leakage, seismic background, and environmental conditions.
\item \textbf{Control and Data Acquisition (CDAQ)}: Centralized timing,
phase locking, data collection, and system monitoring.
\end{enumerate}

\subsection{Operating Principle}

The energy transmission chain proceeds as follows:

\begin{enumerate}
\item Grid AC power $\rightarrow$ RF power at 1.3~GHz via klystrons
\item RF power $\rightarrow$ stored EM energy in SRF cavities
(TE$_{011}$ + TM$_{010}$ superposition)
\item EM-to-$\psi$ conversion via parametric pumping ($2\omega_{\rm RF}
= \Omega_\psi = 2.6$~GHz driven channel, plus $\omega_{\rm RF} =
\Omega_\psi$ parametric channel)
\item $\psi$-mode propagation along the guided channel (1~km)
\item $\psi$-to-EM conversion at receiver cavity array
\item RF rectification $\rightarrow$ DC power output
\end{enumerate}

%=============================================================================
\section{Transmitter Station Design}
%=============================================================================

\subsection{RF Power System}

\subsubsection{Klystron Array}

\begin{table}[H]
\caption{RF power system specifications.}
\begin{tabular}{lll}
\toprule
Parameter & Value & Notes \\
\midrule
Operating frequency & 1.3~GHz & L-band, ILC standard \\
Number of klystrons & 10 & 1 per cavity row \\
Power per klystron & 10~kW CW & Derated from 65~kW pulsed \\
Total RF power & 100~kW & \\
Klystron efficiency & 65\% & Multibeam klystron \\
Grid power for RF & 154~kW & \\
Modulation capability & $m = 0.1$ at $2\omega$ & Phase-locked modulator \\
Phase stability & $< 0.1^\circ$ rms & Cavity-referenced PLL \\
\bottomrule
\end{tabular}
\end{table}

\subsubsection{RF Distribution Network}

Each klystron feeds a row of 10 cavities through a waveguide
distribution manifold with:
\begin{itemize}
\item WR-650 rectangular waveguide (1.3~GHz fundamental mode)
\item Directional couplers for power monitoring
\item Phase shifters (mechanical + electronic) for array phasing
\item Circulator/load protection on each cavity input
\end{itemize}

\subsection{SRF Cavity Array}

\subsubsection{Cavity Design}

\begin{table}[H]
\caption{SRF cavity specifications.}
\begin{tabular}{lll}
\toprule
Parameter & Value & Notes \\
\midrule
Cavity type & TESLA-style elliptical & 9-cell, Nb \\
Fundamental frequency & 1.3~GHz & TM$_{010}$ passband \\
Operating gradient & 10~MV/m & Derated for CW \\
$Q_0$ at operating gradient & $2 \times 10^{10}$ & Nitrogen-doped Nb \\
Stored energy per cavity & 8~J & At 10~MV/m \\
Operating temperature & 2.0~K & Superfluid He \\
Dynamic heat load & 5~W per cavity & At $Q_0 = 2 \times 10^{10}$ \\
TE$_{011}$ mode & Enabled & Dual-mode coupler \\
\bottomrule
\end{tabular}
\end{table}

\subsubsection{TE+TM Superposition}

The critical innovation for $\psi$-mode generation is operating each
cavity in a simultaneous TE$_{011}$ + TM$_{010}$ superposition. This
breaks the geometry transparency (Appendix~R, DFD Review) that
suppresses the driven overlap integral $G$ in pure eigenmodes.

\paragraph{Mode superposition implementation.}
\begin{itemize}
\item Primary coupler: Excites TM$_{010}$ passband (standard)
\item Secondary coupler: Excites TE$_{011}$ mode through orthogonal port
\item Phase relationship: $\phi_{\rm TE-TM}$ controlled to $< 1^\circ$
\item Relative amplitude: $|a_{\rm TE}|/|a_{\rm TM}| = 0.3$--1.0 (tunable)
\end{itemize}

The restored overlap integral with TE+TM superposition:
\begin{equation}
G = u(z_0)\,e^{-2\psi_0}\,\eta_\times\,U_0\,\cos\phi_{\rm TE-TM},
\end{equation}
where $\eta_\times = 0.1$--1.0 depending on the cavity geometry and
mode amplitude ratio.

\subsubsection{Array Configuration}

The 100 cavities are arranged in a $10 \times 10$ planar array:
\begin{itemize}
\item Array dimensions: $5 \times 5$~m footprint
\item Cavity spacing: 0.5~m center-to-center
\item Orientation: Cavity axes parallel, aligned along the
$\psi$-channel direction ($z$-axis)
\item Phase coherence: All cavities locked to common reference
with $\delta\phi < 1^\circ$ ($< 0.02$~rad)
\end{itemize}

The coherent array gain for $\psi$-mode emission:
\begin{equation}
P_\psi^{\rm array} = N_{\rm cav}^2\,P_\psi^{\rm single}
= 10^4 \times P_\psi^{\rm single}.
\end{equation}

\subsection{Cryogenic System}

\begin{table}[H]
\caption{Cryogenic system specifications.}
\begin{tabular}{lll}
\toprule
Parameter & Value & Notes \\
\midrule
Operating temperature & 2.0~K & Superfluid He \\
Static heat load (100 cavities) & 200~W at 2~K & Thermal radiation, conduction \\
Dynamic heat load (100 cavities) & 500~W at 2~K & RF dissipation \\
Total 2~K heat load & 700~W & \\
Refrigerator capacity & 1~kW at 2~K & With 30\% margin \\
He inventory & 2000~L liquid & \\
Refrigerator input power & 250~kW electrical & COP $\sim 1/250$ at 2~K \\
\bottomrule
\end{tabular}
\end{table}

The cryogenic plant is the largest single power consumer in the
facility, typical of SRF installations. A single Linde or Air Liquide
helium refrigerator provides the 1~kW at 2~K capacity.

\subsection{$\psi$-Mode Launcher}

The cavity array is housed within a $\psi$-mode waveguide ``tube''---a
cylindrical vacuum vessel aligned along the channel axis with:
\begin{itemize}
\item Tube diameter: 1.0~m (area $A_\psi = 0.79$~m$^2$)
\item Tube length: 6~m (spanning the cavity array)
\item Cavity apertures at $\psi$-mode antinodes
\item Total cavity aperture: $A_{\rm cav,tot} = 100 \times 3 \times 10^{-3}$~m$^2$
$= 0.3$~m$^2$
\end{itemize}

The area ratio for the parametric overlap:
\begin{equation}
H \approx \frac{2}{6}\times 1.0 \times \frac{0.3}{0.79} \approx 0.13.
\end{equation}

\subsection{$2\omega$ Modulation System}

The amplitude modulation at $2\omega = 2.6$~GHz is implemented by:
\begin{itemize}
\item Vector modulator on each klystron drive signal
\item Modulation depth: $m = 0.1$ (10\% amplitude modulation)
\item Phase accuracy: $< 0.5^\circ$ at 2.6~GHz relative to reference
\item Bandwidth: 10~MHz around $2\omega$ (for tuning/scanning)
\end{itemize}

\subsection{Transmitter Power Budget}

\begin{table}[H]
\caption{Transmitter station power budget.}
\begin{tabular}{lrr}
\toprule
Subsystem & Power (kW) & \% of total \\
\midrule
Cryogenic plant & 250 & 63\% \\
RF klystrons (grid input) & 154 & 39\% \\
Controls, diagnostics, utilities & 20 & 5\% \\
HVAC, lighting, safety & 10 & 3\% \\
\midrule
\textbf{Total TX station} & \textbf{434} & \\
\bottomrule
\end{tabular}
\end{table}

%=============================================================================
\section{Guided $\psi$-Channel Design}
%=============================================================================

\subsection{Channel Geometry}

The guided $\psi$-channel is a cylindrical waveguide of scalar
refractive index perturbation, maintained along the 1~km transmission
path by 10 pump stations.

\begin{table}[H]
\caption{$\psi$-channel parameters.}
\begin{tabular}{lll}
\toprule
Parameter & Value & Notes \\
\midrule
Channel length & 1000~m & \\
Channel radius $R_{\rm ch}$ & 0.5~m & Guided mode radius \\
Background $\psi$ excess & $\Delta\psi_0 \sim 10^{-12}$ & On-axis \\
$\psi$-mode frequency & 2.6~GHz & $= 2\omega_{\rm RF}$ \\
$\psi$-mode wavelength & $\lambda_\psi = c_s/\Omega_\psi \approx 0.12$~m & For $c_s \approx c$ \\
V-parameter & $V \approx 2\pi R_{\rm ch}\sqrt{2\eta\Delta\psi_0}/\lambda_\psi$ & \\
Pump station spacing & 100~m & 10 stations \\
Channel height above ground & 1.5~m & On support posts \\
\bottomrule
\end{tabular}
\end{table}

\subsection{Pump Station Design}

Each pump station maintains the local $\psi$-gradient profile. A
pump station consists of:

\begin{itemize}
\item \textbf{Cavity cluster}: 4 SRF cavities in a compact cryostat,
operated in TE+TM superposition at 1.3~GHz
\item \textbf{Local cryogenics}: Small pulse-tube cryocooler providing
10~W at 4.2~K (cavities at 4.2~K, lower $Q$ but adequate for pumping)
\item \textbf{RF source}: 4~kW solid-state amplifier
\item \textbf{Power supply}: Grid connection or local solar/battery
\item \textbf{Controls}: Fiber-optic link to CDAQ for timing and phase lock
\end{itemize}

\begin{table}[H]
\caption{Pump station specifications (per station).}
\begin{tabular}{lll}
\toprule
Parameter & Value & Notes \\
\midrule
Number of cavities & 4 & TESLA 1-cell \\
Operating temperature & 4.2~K & Pulse-tube cooled \\
$Q_0$ at 4.2~K & $10^8$ & Adequate for pumping \\
Stored energy & 0.4~J per cavity & \\
RF power input & 4~kW & Solid-state amp \\
Cryocooler power & 3~kW & Sumitomo RDE-418 \\
Total station power & 10~kW & \\
Station footprint & $2 \times 2$~m & Weather-protected enclosure \\
\bottomrule
\end{tabular}
\end{table}

\paragraph{Total channel power.} 10 pump stations $\times$ 10~kW = 100~kW
for channel maintenance. This is the ``infrastructure'' power cost of
the $\psi$-waveguide.

\subsection{Channel Alignment and Stability}

The $\psi$-channel must be maintained in a straight line (or
controlled curve) between TX and RX. Requirements:
\begin{itemize}
\item Pump station positional accuracy: $< 10$~cm (relative to channel axis)
\item Phase stability between stations: $< 5^\circ$ at 2.6~GHz
\item Environmental stability: Temperature compensation for
phase drift ($< 1^\circ$/K)
\item Seismic isolation: Pump stations on passive isolators to
prevent microphonics
\end{itemize}

The pump stations are connected by a fiber-optic timing distribution
network providing:
\begin{itemize}
\item 10~MHz reference signal (rubidium-disciplined)
\item 2.6~GHz phase reference (distributed via fiber)
\item Sub-picosecond timing accuracy over 1~km
\end{itemize}

\subsection{Channel Attenuation Estimate}

For the guided mode with $Q_\psi = 10^3$ (conservative):
\begin{equation}
\alpha_\psi = \frac{2\pi}{\lambda_\psi\,Q_\psi}
= \frac{2\pi}{0.12 \times 10^3} = 0.052\,\text{m}^{-1},
\end{equation}
giving power attenuation over 1~km:
\begin{equation}
e^{-\alpha_\psi \times 1000} = e^{-52} \approx 10^{-23}.
\end{equation}

This is unacceptably large. However, the pump stations regenerate
the signal every 100~m. With pump-assisted propagation, the effective
attenuation is:
\begin{equation}
\text{Loss per 100~m segment} = e^{-\alpha_\psi \times 100}
= e^{-5.2} \approx 0.0055,
\end{equation}
and with regeneration at each pump station:
\begin{equation}
\eta_{\rm channel} = \prod_{i=1}^{10}\eta_{\rm segment}
\times \prod_{i=1}^{10}\eta_{\rm pump},
\end{equation}
where $\eta_{\rm pump}$ is the pump station regeneration efficiency.

For $Q_\psi = 10^6$ (optimistic but physically motivated):
\begin{equation}
\alpha_\psi = 5.2 \times 10^{-5}\,\text{m}^{-1}, \quad
e^{-\alpha_\psi \times 100} = e^{-0.0052} \approx 0.995.
\end{equation}

In this case, each 100~m segment has only 0.5\% loss, and the
total channel loss over 1~km is:
\begin{equation}
\eta_{\rm channel} = (0.995)^{10} \approx 0.95.
\end{equation}

The actual $Q_\psi$ will be determined experimentally in Phase~I and
is a critical input to the Phase~II design.

%=============================================================================
\section{Receiver Station Design}
%=============================================================================

\subsection{$\psi$-Mode Collector}

The receiver mirrors the transmitter: a $10 \times 10$ array of SRF
cavities tuned to the $\psi$-mode frequency, with TE+TM superposition
to enable $\psi$-to-EM back-conversion.

The incoming $\psi$-perturbation parametrically drives EM modes in the
receiver cavities. The driven EM amplitude is:
\begin{equation}
|a_{\rm EM}|^2 \propto |\lambda - 1|^2\,|\delta\psi_{\rm in}|^2\,Q_{\rm rec}^2.
\end{equation}

The high $Q_{\rm rec}$ of the SRF cavities provides built-in
amplification of the weak incoming signal.

\subsection{RF Extraction and Rectification}

Power is extracted from the receiver cavities through output couplers
and fed to a rectification system:
\begin{itemize}
\item Output coupler: Variable $Q_{\rm ext}$ matched to loaded $Q$
\item RF power combiner: 100 cavity outputs combined coherently
\item Rectifier: GaN Schottky diode rectenna array
\item DC-DC converter: Boost converter to 48~V DC bus
\end{itemize}

\begin{table}[H]
\caption{Receiver station specifications.}
\begin{tabular}{lll}
\toprule
Parameter & Value & Notes \\
\midrule
Cavity array & $10 \times 10$ SRF & Matched to TX \\
Operating temperature & 2.0~K & \\
Receiver $Q_0$ & $2 \times 10^{10}$ & Same as TX \\
RF extraction power & TBD & Depends on $|\lambda-1|$ \\
Rectifier efficiency & 85\% & GaN Schottky \\
Cryogenic plant & 1~kW at 2~K & Same as TX \\
Station power (non-signal) & 300~kW & Cryo + utilities \\
\bottomrule
\end{tabular}
\end{table}

\subsection{Signal Verification System}

To distinguish genuine $\psi$-mode energy from systematic artifacts,
the receiver includes:

\begin{enumerate}
\item \textbf{EM shielding}: Double-walled copper Faraday room around
the entire receiver cavity array, with $> 120$~dB isolation at 1.3~GHz
and $> 100$~dB at 2.6~GHz.

\item \textbf{Coherence verification}: The receiver signal must be
phase-coherent with the transmitter $2\omega$ modulation (through the
fiber timing link) at the expected $\psi$-channel group delay
$\tau = L/c_s \approx 3.3~\mu$s.

\item \textbf{Null tests}:
\begin{itemize}
\item TX modulation off: Signal should vanish
\item TX in single mode (no TE+TM): Signal should vanish
(geometry transparency)
\item TX frequency detuned from $\Omega_\psi$: Signal should
drop as Lorentzian
\item Lateral RX displacement: Signal should drop outside
channel radius $R_{\rm ch}$
\item Channel pump stations off: Signal should drop
(no waveguide confinement)
\end{itemize}

\item \textbf{Atom interferometer cross-check}: A cold-atom
interferometer co-located with the receiver cavity provides an
independent $\psi$-gradient measurement.
\end{enumerate}

%=============================================================================
\section{Diagnostic Network}
%=============================================================================

\subsection{$\psi$-Field Sensors}

Distributed along the 1~km channel path at 50~m intervals (21 sensors):

\begin{itemize}
\item \textbf{Type}: Compact SRF cavity ($Q \sim 10^8$) in
standalone cryostat with frequency readout
\item \textbf{Sensitivity}: $\delta\psi \sim 10^{-15}$ per
$\sqrt{\text{Hz}}$ (comparable to cavity-atom comparison experiments)
\item \textbf{Measurement}: Cavity frequency shift
$\delta f/f = -2\delta\psi$ from constitutive relations
\item \textbf{Bandwidth}: DC to 10~GHz
\item \textbf{Readout}: Pound-Drever-Hall lock, frequency counted
against fiber-distributed reference
\end{itemize}

\subsection{EM Background Monitors}

At each pump station and at TX/RX:
\begin{itemize}
\item Broadband EM field probes (10~MHz -- 10~GHz)
\item Spectrum analyser monitoring at $\omega$, $2\omega$, $\Omega_\psi$
\item Cross-correlation with TX modulation signal
\end{itemize}

\subsection{Environmental Sensors}

Along the channel path:
\begin{itemize}
\item Temperature (RTDs at 10~m spacing)
\item Barometric pressure (at each pump station)
\item Seismometers (broadband, at TX, RX, and midpoint)
\item Wind speed and direction (for vibration correlation)
\item Ground water table monitors (density variations affect $\psi$)
\end{itemize}

%=============================================================================
\section{Control and Data Acquisition}
%=============================================================================

\subsection{Timing Distribution}

\begin{itemize}
\item Master oscillator: Rubidium-disciplined crystal oscillator
(short-term stability $< 10^{-12}$ at 1~s)
\item Distribution: Single-mode fiber with active path length
stabilization
\item Reference signals: 10~MHz, 1.3~GHz, 2.6~GHz distributed
to all stations
\item Absolute time: GPS-disciplined for data timestamping
\end{itemize}

\subsection{Phase Locking Architecture}

\begin{enumerate}
\item Master reference $\rightarrow$ 10 klystrons (PLL, $< 0.1^\circ$)
\item Klystron $\rightarrow$ 10 cavities per row (cavity field probe feedback)
\item TX array $\rightarrow$ 10 pump stations (fiber phase link)
\item TX array $\rightarrow$ RX array (fiber phase link for
verification; not used for signal)
\end{enumerate}

\subsection{Data Acquisition}

\begin{itemize}
\item ADC sampling: 5~GS/s, 14-bit at TX and RX cavity probes
\item Slow DAQ: 1~kS/s for temperatures, pressures, alignment
\item Data rate: $\sim$1~TB/day (raw), 10~GB/day (processed)
\item Storage: On-site RAID array + cloud backup
\item Analysis: Real-time matched filter for $\psi$-mode signal
at $\Omega_\psi$, with post-processing pipeline for systematic
studies
\end{itemize}

%=============================================================================
\section{Site Requirements}
%=============================================================================

\subsection{Location Criteria}

\begin{enumerate}
\item \textbf{Terrain}: Flat, stable geology, minimal seismic activity.
Preferred: dry lakebed, flat agricultural land, or decommissioned airfield.
\item \textbf{Length}: $> 1.2$~km clear line between TX and RX
\item \textbf{Electrical power}: 1~MW grid connection (or dedicated
substation)
\item \textbf{Access}: Road access for heavy equipment (cryogenic
deliveries, SRF cavity transport)
\item \textbf{EM environment}: Low RF background; away from
transmitter towers, airports, and industrial EM sources
\item \textbf{Regulatory}: Appropriate zoning for research facility;
not near gravitational wave observatories or precision clock laboratories
\end{enumerate}

\subsection{Candidate Sites}

\begin{itemize}
\item SLAC National Accelerator Laboratory (existing SRF infrastructure,
2~km linear footprint)
\item Fermilab (existing SRF test facilities, open land)
\item Jefferson Lab (SRF expertise, adequate space)
\item ESS (Lund, Sweden; new SRF facility with available land)
\item DESY (Hamburg; XFEL tunnel infrastructure)
\end{itemize}

\subsection{Facility Layout}

\begin{verbatim}
[TX Building]---100m---[PS1]---100m---[PS2]---...---[PS10]---[RX Building]
     |                                                            |
     +---------- Fiber optic timing backbone ---------+-----------+
     |                                                            |
     +---------- Diagnostic sensor cable runs --------+-----------+
     |                                                |
  [Control       [Cryo Plant TX]            [Cryo Plant RX]
   Building]
\end{verbatim}

\begin{table}[H]
\caption{Facility building requirements.}
\begin{tabular}{lll}
\toprule
Building & Size & Contents \\
\midrule
TX Hall & $15 \times 15 \times 5$~m & Cavity array, RF distribution \\
TX Cryo Plant & $10 \times 20 \times 8$~m & He refrigerator, compressors \\
RX Hall & $15 \times 15 \times 5$~m & Cavity array, rectification \\
RX Cryo Plant & $10 \times 20 \times 8$~m & He refrigerator, compressors \\
Control Building & $10 \times 15 \times 4$~m & CDAQ, offices, meeting room \\
Pump Station (x10) & $3 \times 3 \times 3$~m & Each with local cryo, RF \\
\bottomrule
\end{tabular}
\end{table}

%=============================================================================
\section{Safety Systems}
%=============================================================================

\subsection{Radiation Safety}

\begin{itemize}
\item SRF cavities can produce X-rays from field emission at high
gradients. Standard accelerator-grade shielding (concrete + lead)
around cavity arrays.
\item $\psi$-field perturbations at $\delta\psi \sim 10^{-12}$ produce
gravitational acceleration perturbations $\delta a \sim 10^{-5}$~m/s$^2$,
negligible compared to tidal accelerations.
\item No ionizing radiation from the $\psi$-channel itself.
\end{itemize}

\subsection{Cryogenic Safety}

\begin{itemize}
\item Oxygen deficiency hazard (ODH) monitoring in all enclosed spaces
\item Pressure relief on all cryogenic vessels (ASME-rated)
\item Helium leak detection system
\item Emergency ventilation in TX and RX halls
\item Personal protective equipment (PPE) for cryogenic operations
\end{itemize}

\subsection{Electrical Safety}

\begin{itemize}
\item High-voltage RF systems: Interlocked access, warning lights
\item Klystron HV supplies: $\sim$100~kV, standard klystron gallery safety
\item Ground fault protection on all circuits
\item Emergency power-off (EPO) buttons at all stations
\end{itemize}

\subsection{EM Exposure}

\begin{itemize}
\item RF leakage monitoring around TX and RX buildings
\item Personnel exclusion zones during high-power operation
\item Compliance with ICNIRP exposure limits
\end{itemize}

%=============================================================================
\section{Cost Estimate}
%=============================================================================

\begin{table}[H]
\caption{Cost breakdown for km-scale demonstration facility.}
\begin{tabular}{lrr}
\toprule
Item & Quantity & Cost (\$M) \\
\midrule
\multicolumn{3}{l}{\textit{Transmitter Station}} \\
\quad SRF cavities (9-cell TESLA) & 100 & 5.0 \\
\quad Cryomodules & 10 & 2.0 \\
\quad Klystrons + HV supplies & 10 & 1.5 \\
\quad RF distribution & 1 set & 0.5 \\
\quad He refrigerator (1~kW @ 2~K) & 1 & 3.0 \\
\quad TX building + shielding & 1 & 1.0 \\
\midrule
\multicolumn{3}{l}{\textit{Guided Channel}} \\
\quad Pump station SRF cavities (1-cell) & 40 & 0.8 \\
\quad Pulse-tube cryocoolers & 10 & 0.5 \\
\quad Solid-state RF amplifiers & 10 & 0.2 \\
\quad Pump station enclosures & 10 & 0.3 \\
\quad Fiber optic timing network & 1~km & 0.2 \\
\midrule
\multicolumn{3}{l}{\textit{Receiver Station}} \\
\quad SRF cavities (9-cell TESLA) & 100 & 5.0 \\
\quad Cryomodules & 10 & 2.0 \\
\quad He refrigerator (1~kW @ 2~K) & 1 & 3.0 \\
\quad Rectenna array + power electronics & 1 set & 0.3 \\
\quad RX building + shielding & 1 & 1.0 \\
\midrule
\multicolumn{3}{l}{\textit{Diagnostics}} \\
\quad $\psi$-field sensor cavities & 21 & 1.0 \\
\quad Atom interferometer & 1 & 1.5 \\
\quad EM monitors + environmental sensors & 1 set & 0.3 \\
\midrule
\multicolumn{3}{l}{\textit{Control \& Infrastructure}} \\
\quad CDAQ system & 1 & 0.5 \\
\quad Control building & 1 & 0.5 \\
\quad Site preparation + utilities & 1 & 1.0 \\
\quad Electrical infrastructure (1~MW) & 1 & 0.5 \\
\midrule
\multicolumn{3}{l}{\textit{Other}} \\
\quad Engineering \& design (15\%) & --- & 4.5 \\
\quad Installation \& commissioning & --- & 2.0 \\
\quad Contingency (20\%) & --- & 6.0 \\
\midrule
\textbf{Total} & & \textbf{42.6} \\
\bottomrule
\end{tabular}
\end{table}

\subsection{Annual Operating Cost}

\begin{table}[H]
\caption{Estimated annual operating costs.}
\begin{tabular}{lr}
\toprule
Item & Cost (\$M/yr) \\
\midrule
Electricity (1~MW, 50\% duty) & 0.4 \\
Liquid helium / cryogen losses & 0.2 \\
Staff (15 FTE) & 2.5 \\
Maintenance \& spares & 0.5 \\
\midrule
\textbf{Total annual operating} & \textbf{3.6} \\
\bottomrule
\end{tabular}
\end{table}

%=============================================================================
\section{Schedule}
%=============================================================================

\begin{table}[H]
\caption{Project schedule (36 months).}
\begin{tabular}{lll}
\toprule
Phase & Duration & Activities \\
\midrule
Design \& procurement & Months 1--12 &
  Detailed design, cavity fabrication, \\
  & & site preparation, long-lead procurement \\
\midrule
Installation & Months 10--24 &
  Building construction, cryoplant install, \\
  & & cavity installation, RF commissioning \\
\midrule
Cold commissioning & Months 22--28 &
  Cooldown, cavity qualification, \\
  & & single-cavity $\psi$-tests (Phase I repeat) \\
\midrule
Array commissioning & Months 26--32 &
  Phase-lock 100-cavity array, \\
  & & pump station commissioning \\
\midrule
$\psi$-channel tests & Months 30--36 &
  Channel formation verification, \\
  & & end-to-end power transmission, \\
  & & systematic studies \\
\bottomrule
\end{tabular}
\end{table}

%=============================================================================
\section{Success Criteria and Milestones}
%=============================================================================

\subsection{Go/No-Go Decision Points}

\begin{enumerate}
\item \textbf{Month 6 --- Phase I validation}: Repeat the laboratory
$\psi$-mode generation experiment with a single cavity pair. Confirm
$|\lambda - 1| > 10^{-12}$ (minimum for practical demonstration).
\textbf{If null: reassess programme.}

\item \textbf{Month 18 --- Single-cavity channel test}: Demonstrate
$\psi$-mode propagation over 10~m with a single TX cavity and single
RX cavity, no waveguide. Measure $\psi$-mode velocity $c_s$ and
damping rate $\gamma_\psi$. \textbf{If $Q_\psi < 100$: redesign channel.}

\item \textbf{Month 28 --- Array coherence}: Demonstrate $N^2$ power
scaling with $N = 10$ cavities. Confirm phase-locking stability over
1~hour. \textbf{If coherence time $< 1$~s: redesign phase locking.}

\item \textbf{Month 34 --- Channel formation}: Verify guided-mode
confinement using $\psi$-sensor array. Signal should be confined to
$R_{\rm ch}$ and constant along $z$. \textbf{If no confinement: channel
concept requires revision.}

\item \textbf{Month 36 --- Power transmission}: Measure end-to-end
power delivery with all null tests passed. Target: $> 0.1\%$ efficiency
or $> 1$~W received at $> 100$~kW transmitted.
\end{enumerate}

\subsection{Primary Success Metrics}

\begin{table}[H]
\caption{Success criteria for the demonstration.}
\begin{tabular}{lll}
\toprule
Metric & Threshold & Goal \\
\midrule
$\psi$-mode detected at RX & Yes/No & Yes \\
Signal SNR & $> 3$ & $> 100$ \\
End-to-end efficiency & $> 10^{-4}$ & $> 10^{-2}$ \\
Channel confinement & Qualitative & $< 2R_{\rm ch}$ \\
Null tests passed & 3/5 & 5/5 \\
$|\lambda - 1|$ measured & Yes & $\pm 10\%$ \\
$Q_\psi$ measured & Yes & $\pm 20\%$ \\
$c_s$ measured & Yes & $\pm 1\%$ \\
\bottomrule
\end{tabular}
\end{table}

%=============================================================================
\section{Risk Assessment}
%=============================================================================

\begin{table}[H]
\caption{Risk matrix for the demonstration programme.}
\begin{tabular}{p{4.5cm}ccp{5cm}}
\toprule
Risk & Likelihood & Impact & Mitigation \\
\midrule
$|\lambda - 1| < 10^{-14}$
(no detectable coupling) & Medium & Critical &
Phase I go/no-go at Month 6; minimal investment before confirmation \\
\midrule
$Q_\psi < 100$
(high $\psi$-damping) & Medium & High &
Closer pump station spacing; regenerative relay approach \\
\midrule
Array phase coherence
insufficient & Low & High &
Proven SRF technology; digital LLRF feedback \\
\midrule
Systematic EM leakage
mimics signal & Medium & Medium &
120~dB shielding; 5 independent null tests;
atom interferometer cross-check \\
\midrule
SRF cavity performance
degradation in CW & Low & Medium &
N-doped Nb; proven at LCLS-II \\
\midrule
Cost overrun ($> 50\%$) & Medium & Medium &
20\% contingency; phased procurement \\
\midrule
Environmental factors
(seismic, thermal) & Low & Low &
Site selection; isolation; compensation \\
\bottomrule
\end{tabular}
\end{table}

%=============================================================================
\section{Path to Commercial Deployment}
%=============================================================================

\subsection{Phase III: 100~km, 100~MW System}

Assuming Phase~II success, the path to a commercial-scale system involves:

\begin{enumerate}
\item \textbf{Cavity cost reduction}: From \$50k to $\sim$\$1k per
cavity through mass production (cf.\ automotive SRF: crabbing cavities
for HL-LHC produced at $\sim$\$20k each in small batches).

\item \textbf{Cryocooler integration}: Replace large He plants with
distributed cryocoolers at each station, eliminating central cryogenics
and He logistics.

\item \textbf{Higher temperature operation}: Nb$_3$Sn or MgB$_2$
cavities operating at 4.2~K or above, dramatically reducing cryogenic
cost.

\item \textbf{Array scaling}: $10^4$--$10^6$ cavity arrays, potentially
using thin-film SRF on copper substrates for low-cost mass production.

\item \textbf{$\psi$-channel optimization}: Engineered background
$\psi$-profiles using permanent density structures (underground mass
distributions) to reduce active pumping requirements.
\end{enumerate}

\subsection{Target Commercial Parameters}

\begin{table}[H]
\caption{Commercial system targets (Phase III+).}
\begin{tabular}{lll}
\toprule
Parameter & Phase II Demo & Commercial Target \\
\midrule
Distance & 1~km & 100--1000~km \\
Power & $\sim$1~W received & 10--100~MW \\
Efficiency & 0.1\% & $> 50\%$ \\
Capital cost & \$43M & \$1--3M per km \\
Operating cost & \$3.6M/yr & \$50k/yr per km \\
Cavity count & 200 (TX+RX) & $10^4$--$10^6$ \\
\bottomrule
\end{tabular}
\end{table}

%=============================================================================
\section{Intellectual Property Coverage}
%=============================================================================

This engineering design falls under the scope of U.S. Provisional
Patent Application 63/865,283 (filed August 16, 2025), which covers:

\begin{enumerate}
\item Systems and methods for energy transmission via $\psi$-field modulation
\item Transmitter apparatus using EM cavity arrays to generate $\psi$-perturbations
\item Guided $\psi$-channel maintained by distributed pump stations
\item Receiver apparatus converting $\psi$-mode energy to electrical power
\item Applications including wireless power distribution, space power downlink,
subsurface energy delivery, and directed-energy propulsion
\end{enumerate}

%=============================================================================
\section{Conclusions}
%=============================================================================

This engineering design presents a technically feasible path to demonstrating
km-scale wireless power transmission via the DFD $\psi$-channel mechanism.
The design leverages mature SRF cavity technology (from particle accelerator
programmes), proven cryogenic infrastructure, and established RF engineering
to create an experimental facility capable of:

\begin{enumerate}
\item Definitively confirming or refuting $\psi$-mode generation and propagation
\item Measuring the critical parameters $|\lambda - 1|$, $Q_\psi$, and $c_s$
\item Demonstrating guided $\psi$-channel energy transmission over 1~km
\item Establishing the efficiency baseline for scaling to commercial systems
\end{enumerate}

The total estimated cost of \$43M over 36 months is comparable to a
mid-scale accelerator test facility and is justified by the transformative
potential of the technology. The staged go/no-go decision structure ensures
that investment is minimized if early results are negative.

The critical path item is the Phase~I validation of $|\lambda - 1| \neq 0$
at detectable levels. This single measurement, achievable with existing
apparatus in an afternoon, determines the viability of the entire programme.

\end{document}
